\documentclass[10pt]{article}
\usepackage[top=0.74in,bottom=0.76in,left=0.9in,right=0.9in]{geometry}
\usepackage[cmex10]{amsmath}
\usepackage{amsthm}
\usepackage[hidelinks]{hyperref}
\DeclareFontShape{OMX}{cmex}{m}{n}{<->cmex10}{}
\DeclareMathSizes{10}{10}{7}{7}
\DeclareMathSizes{9}{10}{7}{7}
\newtheorem{theorem}{Theorem}
\title{Small incomparability components after one deletion}
\author{Achita Anantachina}
\date{26 November 2026}
\begin{document}
\maketitle

\begin{abstract}
A hypothetical counterexample to the one-third--two-thirds conjecture
cannot have a small component in the incomparability graph after deleting
any one point. We give an elementary proof for components of order two and
three and an exact finite check that treats larger components at arbitrary
ambient order. We also search for counterexamples by inserting one or two
points into published low-balance posets.
\end{abstract}

For a finite poset \(Q\), a balanced pair consists of incomparable
\(u,v\) for which the probability of \(u\) preceding \(v\) in a uniform
linear extension lies in \([1/3,2/3]\). The conjecture asks for such a
pair in every non-chain. See \cite{acpp} for recent general bounds
and \cite{gupta} for the order-14 census.

\begin{theorem}\label{thm:small}
Let \(Q\) be a finite non-chain poset and \(x\in Q\). If \(Q-x\) is a chain,
or if its incomparability graph has a component of order two or three,
then \(Q\) has a balanced pair.
\end{theorem}

\begin{proof}
If \(Q-x\) is a chain, the elements incomparable with \(x\) form a
nonempty interval. Inserting \(x\) in a gap of that interval gives all
extensions. A middle element of the interval gives a balanced pair.

Otherwise let \(B\) be a component of order two or three after deletion.
The components of any incomparability graph form an ordinal sum in the
poset. Indeed, an outside vertex compares in the same direction to both
ends of every incomparability edge, by transitivity. This fixes the
direction between entire components. Every extension is then a
concatenation of component extensions. We may restrict to the component
of the incomparability graph of \(Q\) containing \(B\). If it omits \(x\), it
equals \(B\) because all old incomparability edges are unchanged.
An antichain of order two or three has a pair with probability \(1/2\).
The only other connected possibility of order three is \(a<b\) with
\(c\) incomparable to both, where \(b\) precedes \(c\) with probability
\(1/3\).

Suppose the component contains \(x\). If there are at least two old
components in it, an element below \(x\) can only be in the first, and
an element above \(x\) can only be in the last. For example a predecessor
of \(x\) in any later component would make every element of the first
component comparable to \(x\), disconnecting it. The upper statement
follows by duality. If \(B\) is interior, \(x\) is incomparable to all its
points. For every fixed choice of the other component extensions,
the number of admissible gaps for \(x\) does not depend on the order
within \(B\). Each extension of \(B\) thus has equal weight, giving the
same balanced pair as in \(B\).

At an end we treat the first component, using duality for the last.
Here \(x\) has no successor in \(B\). If \(B\) is the only old component
and all comparisons with \(x\) point one way, choose the dual if needed
so that \(x\) has no successor in \(B\).
For each extension of \(B\), let \(\ell\) be the last position of an
element below \(x\), or zero if there is none. There are \(|B|-\ell\)
admissible gaps before the end of \(B\). For each fixed extension of the
rest, the number of gaps at or after the end is independent of the order
within \(B\) and is at least one. After averaging over the rest, the
weights of the orders of \(B\) are proportional to \(r+|B|-\ell\) for
a common \(r\ge1\).

For a two-point antichain the two weights differ by at most one, and
both are at least one. Their normalized values are in \([1/3,2/3]\).
For a three-point antichain, two points have the same relation to \(x\).
They can be swapped by an automorphism, giving probability \(1/2\).
For \(B=\{a<b,c\}\), use the orders \(abc,acb,cab\). As the ideal of
predecessors varies through \(\emptyset,\{a\},\{c\},\{a,c\},\{a,b\},B\),
the three \(\ell\)-values are respectively
\[
(0,0,0),\ (1,1,2),\ (3,2,1),\ (3,2,2),\ (2,3,3),\ (3,3,3).
\]
Thus \(abc\) or \(cab\) has maximum weight. All three weights are at
least one and differ by at most two, so the largest normalized weight
is between \(1/3\) and \(3/5\). The events \(abc\) and \(cab\) are exactly
\(b\) before \(c\) and \(c\) before \(a\), respectively.

Finally, if \(B\) is the only old component, both a predecessor and
a successor of \(x\) can occur only in the three-point case with
\(a<x<b\). Then \(c\) is isolated from the chain, and \(c,x\) is balanced.
\end{proof}

In particular, all one-point extensions of ordinal sums of singleton
and three-point equality blocks satisfy the conjecture whenever they
are non-chains. Theorem~\ref{thm:small} applies even when other old
components are arbitrary. A counterexample must, after every one-point
deletion, have at least one incomparability component of order at least
four and none of order two or three.

\section*{Checking larger components}
The boundary weights above reduce larger components to a finite problem
with one real parameter. For an endpoint component \(B\) and its ideal
\(D\) of predecessors of \(x\), set \(o(\pi)=|B|-\ell(\pi)\) for each
linear extension \(\pi\) of \(B\). Write \(m\) for the number of
extensions and \(h=\sum_\pi o(\pi)\). For any incomparable pair \(u,v\)
write \(c\) for the number of extensions with \(u\) before \(v\) and \(d\)
for the sum of their offsets. The probability of this event is
\[
\frac{rc+d}{rm+h},\qquad r\ge1.
\]
The condition of being in \([1/3,2/3]\) gives two affine inequalities
with integer coefficients. Intersect their closed solution intervals
with \([1,\infty)\), and take the union over incomparable pairs.
This gives an exact test for coverage of every admissible parameter.
For example, with \(a<b<c\), isolated \(d\), and \(D=\{d\}\), the pair
\(d,a\) is balanced for \(1\le r\le3\), while \(d,b\) is balanced for
\(r\ge3/2\). Both intervals are needed to cover the whole range.

The enumeration uses natural labellings, in which every relation points
from a smaller to a larger label. Any poset has one, and a new largest
label can have any ideal as its strict lower set. We generate all such
orders up to the tested size, keep those with connected incomparability
graph, and check every ideal \(D\).

Through order seven every interval union covered \([1,\infty)\). The
order-seven check tested 86,278 connected naturally labelled orders and
2,589,300 choices of an order and an ideal. This extends
Theorem~\ref{thm:small} to components of orders four through seven.
For an interior component the empty ideal gives its uniform law. If
the component is the entire old order after restriction and \(x\) has both
predecessors and successors, the restricted poset has at most eight
points and is covered by the order-14 census \cite{gupta}.

\section*{Exact checks near low-balance seeds}
We use the 42 labelled representatives of the distinct low-tail
values in Gupta's archive and the 68 stored equality representatives.
For an inserted vertex its strict lower set is an ideal \(D\), its
upper set is a filter \(U\), and every point of \(D\) must be strictly
below every point of \(U\). These conditions are necessary and
sufficient. Every legal insertion appears once for each labelled seed.

The exact counting algorithm traverses the ideals of the resulting
poset. If an edge of the ideal graph inserts \(v\) into ideal \(I\), the
product of the prefix count at \(I\) and suffix count at \(I+v\) is the
number of extensions using this edge. Add this product to the
precedence count of \(u,v\) for each \(u\in I\). Every extension is
counted exactly once for each ordered precedence it contains.

For one insertion we checked 27,371 labelled cases. Among the 7,315
cases with connected incomparability graph, the minimum balance was
\(889/2483\). For two insertions into the 42 tail representatives,
6,899,699 labelled cases gave 4,281,043 with connected incomparability
graph, all of which were checked. Their minimum balance was
\(269/767\). Two insertions into the 68 equality representatives gave
3,563,408 labelled cases, of which 336,004 were connected and checked.
Their minimum was \(63/160\). There were no counterexamples in the
checked cases. The coverage is limited to insertions into these seeds.

\paragraph{Acknowledgement.}
The computation and exposition used AI assistance.

\begin{thebibliography}{9}
\bibitem{acpp} M. Aires, S. H. Chan, I. Pak, and G. Panova.
Breaking the infinite barrier in the one-third--two-thirds conjecture.
\href{https://arxiv.org/abs/2609.30888v2}{arXiv:2609.30888v2} (2026).
\bibitem{gupta} A. Gupta. Balance constants, majority cycles, and the
Gold Partition Conjecture through fourteen elements.
\href{https://arxiv.org/abs/2607.23926v2}{arXiv:2607.23926v2} (2026).
\end{thebibliography}
\end{document}
